\documentclass[11pt,twoside]{article}
\newcommand{\thang}[1]{{\color{blue}\small thang:#1}}
\newcommand{\bin}[1]{\langle #1 \rangle_2}
\usepackage[table]{xcolor}
\newcommand{\immune}{\cellcolor{gray} -1 }
\newcommand{\negative}{\cellcolor{green} 0 }
\newcommand{\positive}{\cellcolor{red} 1}

\usepackage{amssymb}
\usepackage{amsmath}
\usepackage{latexsym}
\usepackage{color}
\usepackage{graphics}
\usepackage{xspace}
\usepackage{xlop}
% Commands for special characters
\newcommand{\mybackslash}{\char'134}
\newcommand{\myleftbracket}{\char'173}
\newcommand{\myrightbracket}{\char'175}
\newcommand{\mypercent}{\char'045}
\newcommand{\myunderscore}{\char'137}
\usepackage{array}
\setlength\extrarowheight{2pt}
% The 'ifthen' package supports Boolean flags
\usepackage{ifthen}
% The 'solutions' flag determines whether this is the original handout
%    or a solution
\newboolean{solutions}
\setboolean{solutions}{False}  % Default is original handout

% Uncomment the next line before starting on the solutions
% \setboolean{solutions}{True}

\newcommand{\coursenumber}{CS 4104}
\newcommand{\coursetitle}{Data and Algorithm Analysis}
\newcommand{\docdate}{October 17, 2022}
\newcommand{\duedate}{October 21, 2022}
\newcommand{\homeworknumber}{3}

% The document title depends on whether these are solutions
\ifthenelse{\boolean{solutions}}{% Then
\newcommand{\doctitle}{Solutions to Homework Assignment \homeworknumber}
}{% Else
\newcommand{\doctitle}{Midterm}
}

% The document date depends on whether these are solutions
\ifthenelse{\boolean{solutions}}{% Then
\renewcommand{\docdate}{\duedate}
}{% Else
}

% If you are the author, so put your name here
\renewcommand{\author}{Lenwood S. Heath}

\pagestyle{myheadings}
\markboth{\hfill\doctitle\hfill\docdate}{\docdate\hfill\doctitle\hfill}

\addtolength{\textwidth}{1.00in}
\addtolength{\textheight}{1.00in}
\addtolength{\evensidemargin}{-1.00in}
\addtolength{\oddsidemargin}{-0.00in}
\addtolength{\topmargin}{-.50in}
\setlength{\footskip}{0pt}

\newcommand{\polyreduce}{\leq_{\mathrm{P}}}

\newcounter{problem}
\setcounter{problem}{0}
\newcommand{\problem}[1]{%
\refstepcounter{problem}\noindent\textbf{[#1] \arabic{problem}.}}

\newcommand{\solution}{\bigskip\hrule\bigskip}
\newcommand{\problembreak}{\bigskip\hrule\bigskip}

\renewcommand{\theenumi}{\textbf{\Alph{enumi}}}
\renewcommand{\theenumii}{\textbf{\roman{enumii}}}

\newcommand{\emptysequence}{\Lambda}
\newcommand{\emptystring}{\lambda}

% Pseudocode comment symbol
\newcommand{\comment}{\textbf{//}\ \ }

% Pseudocode line numbering
\newcounter{pseudocode}
\newcommand{\firstline}{\setcounter{pseudocode}{0}\linenumber}
\newcommand{\linenumber}{\refstepcounter{pseudocode}\thepseudocode}
\newcommand{\pseudoindent}{\hspace*{26pt}}

\newcommand{\nil}{\mbox{\textsc{nil}}}

% Mathematical symbols
\newcommand{\grid}{\mathcal{G}}  % Grid graph
\newcommand{\integers}{\mathbb{Z}}  % Integers
\newcommand{\naturals}{\mathbb{N}}  % Natural numbers
\newcommand{\reals}{\mathbb{R}}  % Real numbers

% Probability
\newcommand{\expect}[1]{\mathbf{E}\left[#1\right]}
\newcommand{\prob}[1]{\mathrm{Pr}\left[#1\right]}
\newcommand{\var}[1]{\mathrm{Var}\left[#1\right]}

% Logic
\newcommand{\NOT}[1]{\neg #1}
\newcommand{\AND}{\wedge}
\newcommand{\OR}{\vee}
\newcommand{\clause}[1]{\left(#1\right)}

\newlength{\problemoffset}
\setlength{\problemoffset}{0.5in}

% Decision problem macro
% A command for formatting decision problems a la Garey and Johnson
\newcommand{\decision}[3]{%     \decision{NAME}{INSTANCE}{QUESTION}
\begin{list}{}{
\setlength{\leftmargin}{\problemoffset}
\setlength{\rightmargin}{\problemoffset}
\setlength{\parsep}{0pt}
\setlength{\itemsep}{2pt}
\setlength{\topsep}{\itemsep}
\setlength{\partopsep}{\itemsep}
}
\item
{\textsc{#1}}
\item
{INSTANCE: #2}
\item
{QUESTION: #3}
\end{list}
}

% Optimization problem macro
\newcommand{\optimization}[3]{%  \optimization{NAME}{INSTANCE}{SOLUTION}
\begin{list}{}{
\setlength{\leftmargin}{\problemoffset}
\setlength{\rightmargin}{\problemoffset}
\setlength{\parsep}{0pt}
\setlength{\itemsep}{2pt}
\setlength{\topsep}{\itemsep}
\setlength{\partopsep}{\itemsep}
}
\item
{\rule{0pt}{14pt}\textsc{#1}}
\item
{INSTANCE: #2}
\item
{SOLUTION: #3}
\end{list}
}

\newcommand{\reaches}{\leadsto}

% Table layout
\newcommand{\toprule}{\rule[11pt]{0pt}{2pt}}
\newcommand{\bottomrule}{\rule[-6pt]{0pt}{0pt}}
\newenvironment{raggedpars}[1][2.0in]{%
\begin{minipage}[t]{#1}\raggedright\toprule}%
{\bottomrule\end{minipage}}

% Dynamic programming macros
\newlength{\arrowwidth}
\setbox3=\hbox{$\nwarrow$}
\setlength{\arrowwidth}{\wd3}
\newcommand{\optnwarrow}[1]{\if1#1\nwarrow%
\else\rule{\arrowwidth}{0pt}\fi}
\newcommand{\optuparrow}[1]{\if1#1\uparrow%
\else\rule{\arrowwidth}{0pt}\fi}
\newcommand{\optleftarrow}[1]{\if1#1\leftarrow%
\else\rule{\arrowwidth}{0pt}\fi}
% Use \tablebox to specify any combination of arrows, plus the value
\newcommand{\tablebox}[4]{%
\setlength{\arraycolsep}{0.0pt}%
\begin{array}{cc}%
\optnwarrow{#1} & \optuparrow{#2} \\%
\optleftarrow{#3} & #4%
\end{array}%
}
% Use \tableboxred to specify any combination of arrows, plus the value
% The value will be red; arrows are made red if 2 used instead of 1
\newcommand{\optnwarrowred}[1]{\if1#1\nwarrow%
\else\if2#1\textcolor{red}{\nwarrow}\else\rule{\arrowwidth}{0pt}\fi\fi}
\newcommand{\optuparrowred}[1]{\if1#1\uparrow%
\else\if2#1\textcolor{red}{\uparrow}\else\rule{\arrowwidth}{0pt}\fi\fi}
\newcommand{\optleftarrowred}[1]{\if1#1\leftarrow%
\else\if2#1\textcolor{red}{\leftarrow}\else\rule{\arrowwidth}{0pt}\fi\fi}
\newcommand{\tableboxred}[4]{%
\setlength{\arraycolsep}{0.0pt}%
\begin{array}{cc}%
\optnwarrowred{#1} &%
\optuparrowred{#2} \\%
\optleftarrowred{#3} &%
\textcolor{red}{#4}%
\end{array}%
}

% Allow really sloppy paragraphs that do not generate overfull or
%    underfull hbox's
\newenvironment{SLOPPY}{\begin{sloppypar}\hbadness=10000}{\end{sloppypar}}

% Definitions for this homework
\newcommand{\extent}[1]{\mathrm{extent}(#1)}
\newcommand{\opt}[2]{\mbox{\textsc{opt}}(#1,#2)}
\newcommand{\gap}{\mbox{\texttt{-}}}
\newcommand{\rewriterule}[2]{#1\to #2}
\newcommand{\rewrites}[2][]{\mathop{\Longrightarrow}\limits_{#1}^{#2}}
\newcommand{\reduces}{\leq}
\newcommand{\polyreduces}{\leq_P}
\newcommand{\mathsc}[1]{\mbox{\normalfont\textsc{#1}}}
\newcommand{\NP}{\mathcal{NP}}
\renewcommand{\P}{\mathcal{P}}
\newcommand{\describes}{\vdash}

\begin{document}

\thispagestyle{empty}

\begin{center}
\begin{tabular}{lcr}
\multicolumn{3}{c}{\Large\textbf{\coursenumber}}
\\[0.04in]
\multicolumn{3}{c}{\Large\textbf{\doctitle}}
% If these are solutions, then include the author's (student's) name
\ifthenelse{\boolean{solutions}}{% Then
\\[0.04in]\multicolumn{3}{c}{\large\textbf{\author}}
}{} % Else omit
% If these are solutions, then the date is the due date
\ifthenelse{\boolean{solutions}}{% Then
\\[0.10in]\multicolumn{3}{c}{\duedate}
}{% Else, put given and due dates
\\[0.10in]
\textbf{Given:} \docdate
& \hspace*{1.0in} &
\textbf{Due:} \duedate
}
\end{tabular}
\end{center}

% If these are solutions, then we do not include the directions
\ifthenelse{\boolean{solutions}}{} % No directions
{
% Original document includes directions
\begingroup % This allows an argument that contains multiple paragraphs
\paragraph{General directions.}
The point value of each problem is shown in [ ].
Each solution must include all details and
an explanation of why the given solution is correct.
\textbf{\textcolor{red}{In particular,
write complete sentences.
A correct answer without an explanation is worth no credit.}}
The completed assignment must be submitted on Canvas as a PDF by 5:00 PM
on \duedate.
\textbf{No late homework will be accepted.}

\paragraph{Digital preparation of your solutions is mandatory.
This includes digital preparation of any drawings.}
Use of \LaTeX\ is required.
Also,
\textbf{please include your name.}

\paragraph{Use of \LaTeX\ (required).}
\begin{itemize}
\item Retrieve this \LaTeX\ source file,
named
\texttt{midterm.tex},
from the course web site.
\item Rename the file
\textit{$<$Your VT PID$>$}\texttt{{\myunderscore}solvemidterm.tex},
For example,
for the instructor,
the file name would be
\texttt{heath{\myunderscore}solvemidterm.tex}.

\item
Use a \textbf{text editor}
(such as \texttt{vi}, \texttt{emacs}, or \texttt{pico})
to accomplish the next three steps.
Alternately,
use Overleaf as your \LaTeX\ platform.

\item
Uncomment the line

\texttt{{\mypercent} 
{\mybackslash}setboolean{\myleftbracket}solutions{\myrightbracket}%
{\myleftbracket}True{\myrightbracket}}

in the document preamble by deleting the \texttt{\mypercent}.

\item
Find the line

\texttt{{\mybackslash}renewcommand%
{\myleftbracket}{\mybackslash}author{\myrightbracket}%
{\myleftbracket}Lenwood S. Heath{\myrightbracket}}

and replace the instructor's name with your name.

\item
Enter your solutions where you find
the \LaTeX\ comments

\texttt{{\mypercent} PUT YOUR SOLUTION HERE}

\item
Generate a PDF and turn it in on Canvas by 5:00 PM on \duedate.
\end{itemize}
\endgroup

\problembreak

\newpage

}

\problem{50}
In this game, which we will call the coins-in-a-line game, an even number $n$ of  coins are placed in a line,
with each coin having a known value. Two players, who we will call Reema and Timothy, take turns 
removing one of the coins from either end of the remaining line of coins.
That is, when it is a player’s turn, they remove the coin at the left or right end of the line of coins and add that coin to their collection.
The player who removes a set of coins with a larger total value than the other player wins,
where we assume that both Reema and Timothy know the value of each coin.
Say that Reema always goes first.

{\bfseries

\begin{enumerate}

\item
Provide notation to identify the instances and subinstances of the coins-in-a-line game.

\item
Develop a dynamic programming recurrence
to determine who wins the game and how to play optimally.

\item
Apply your recurrence to an instance of size $n=6$ of your choosing.

\end{enumerate}

}

\solution

% PUT YOUR SOLUTION HERE

\problembreak

\problem{50}
Here is the \textsc{Shortest Cycle} problem.
\optimization{Shortest Cycle}
{$G=(V,E)$, an undirected connected graph.}
{A cycle in $G$ of shortest length or $\nil$ if there is no cycle in $G$.}
Note first that a shortest cycle,
if one exists,
can have length as large as $|V|$.

\textbf{Decide on an algorithmic approach to solve {\normalfont\textsc{Shortest Cycle}}
and explain it in a few sentences.
Give proper pseudocode for your resulting algorithm.
Analyze the worst-case time complexity of your algorithm.}

\solution

% PUT YOUR SOLUTION HERE

\problembreak

\problem{50}
Consider this recurrence for a time complexity $T(n)$:
\begin{eqnarray*}
T(n)
& = &
\left\{
\begin{array}{ll}
T(n/3) +T(n/4)+T(n/5) + c_1 n & n>1;
\\
c_2 & n=1.
\end{array}
\right.
\end{eqnarray*}

\textbf{Solve the above recurrence in the $\Theta$ sense.
Explain carefully the method you use to solve it.
You are not required to draw a recursion tree,
but you may want to describe one.}

\solution

% PUT YOUR SOLUTION HERE

\problembreak

\clearpage

\problem{50}
Let $A$ be an $m\times n$ matrix with every entry from the set $\{-1,0,1\}$.
The matrix represents a region in which a virus can spread,
with each entry a cell in that region.
A -1 cell is immune, that is, cannot be infected.
A 1 cell is positive, that is, infected.
A 0 cell is negative but is susceptible to being infected by a neighbor.
In one time step,
the virus can spread to all of its neighbors
in the four directions left, right, up, and down.

\textbf{Design an algorithm that,
given an initial matrix $A_0$,
returns the number of time steps it takes for the virus to spread in the matrix (region)
and then stop. Clearly indicate the complexity of your designed algorithm.}

\paragraph{Example.}

\begin{equation*}
A_0 = 
\begin{array}{ccccc}
\immune & \negative & \immune & \negative &  \immune   \\
\positive & \immune & \negative & \negative &  \immune   \\
\immune & \negative & \positive & \negative &  \negative   \\
\negative & \immune & \immune  & \negative &  \positive   \\
\immune & \negative & \negative & \negative & \immune   \\
\end{array}
\end{equation*}

\begin{equation*}
A_1 = 
\begin{array}{ccccc}
\immune & \negative & \immune & \negative &  \immune   \\
\positive & \immune & \positive & \negative &  \immune   \\
\immune & \positive & \positive & \positive &  \positive   \\
\negative & \immune & \immune  & \positive &  \positive   \\
\immune & \negative & \negative & \negative & \immune   \\
\end{array}
\end{equation*}


\begin{equation*}
A_2 = 
\begin{array}{ccccc}
\immune & \negative & \immune & \negative &  \immune   \\
\positive & \immune & \positive & \positive &  \immune   \\
\immune & \positive & \positive & \positive &  \positive   \\
\negative & \immune & \immune  & \positive &  \positive   \\
\immune & \negative & \negative & \positive & \immune   \\
\end{array}
\end{equation*}

\begin{equation*}
A_3 = 
\begin{array}{ccccc}
\immune & \negative & \immune & \positive &  \immune   \\
\positive & \immune & \positive & \positive &  \immune   \\
\immune & \positive & \positive & \positive &  \positive   \\
\negative & \immune & \immune  & \positive &  \positive   \\
\immune & \negative & \positive & \positive & \immune   \\
\end{array}
\end{equation*}


\begin{equation*}
A_4 = 
\begin{array}{ccccc}
\immune & \negative & \immune & \positive &  \immune   \\
\positive & \immune & \positive & \positive &  \immune   \\
\immune & \positive & \positive & \positive &  \positive   \\
\negative & \immune & \immune  & \positive &  \positive   \\
\immune & \positive & \positive & \positive & \immune   \\
\end{array}
\end{equation*}

So it requires four time steps
for the virus to spread in the region and stop.

\solution

% PUT YOUR SOLUTION HERE

\problembreak


\end{document}
